\documentstyle[% 


righttag% 


,11pt% 


,amssymb% 


,russian%               remove in Eng. 


,izvru%                 Set izven% in Eng. 


]{amsart} 


 


%  


 


\newcommand{\myfrac}[2]{\frac{\partial #1}{\partial #2}} 


\newcommand{\Tr}{\operatorname{Tr}} 


\newcommand{\im}{\operatorname{Im}} 


\newcommand{\tts}[1]{} 


 


\theoremstyle{plain} 


\theorembodyfont{\it} 


\newtheorem{thms}{\hskip\parindent Теорема}[section] 


\newtheorem{lems}{\hskip\parindent Лемма}[section] 


 


\theoremstyle{definition} 


\theorembodyfont{\rm} 


\newtheorem{cors}{\hskip\parindent Следствие}[section] 


\newtheorem{notes}{\hskip\parindent Замечание}[section] 


\newtheorem{defns}{\hskip\parindent Определение}[section] 


\newtheorem{stats}{\hskip\parindent Утверждение}[section] 


\numberwithin{equation}{section} 


 


%\hoffset=-15mm%                  For Eng. 


  


\setcounter{page}{51} 


\god{1999} 


\nomer{1~(440)} %                        Remove in Eng. 


%\nomer{Vol.\,43, No.\,1}%               Set in Eng. 


%\str{\pageref{begin}--\pageref{end}}%  Set in Eng. 


\udk{517.957} 


\author{\it Л.А.\,Сахнович} 


% NB:   ^^ remove in Eng. 


\title{Матричное нелинейное уравнение Шредингера\\ 


и соответствующая иерархия уравнений} 


 


\thanks{Работа выполнена при частичной финансовой 


поддержке Международного научного Фонда Г.\,Сороса 


(грант UCX\,000).} 


\date{} 


% \pagestyle{myheadings}%         Set in Eng. 


 


\pagestyle{plain} %              Remove in Eng. 


\begin{document} 


 


% \thispagestyle{plain}%          Set in Eng. 


 


\maketitle 


\label{begin} 


 


В статье рассматривается семейство уравнений 


\begin{equation} 


\myfrac{G}{t}-\myfrac{F_N}{x}+[G,F_N]=0;\quad 


N=1,2,3,\dotsc,\tag{0.1} 


\end{equation} 


где $[\cdot,\cdot]$ --- символ коммутанта, т.\,е.~$[G,F_N]=GF_N-F_NG$. 


Матрицы $G(x,t,z)$, $F_N(x,t,z)$ имеют порядок $2m\times 2m$ и 


являются многочленами относительно $z$ 


соответственно степени $1$ и $N$. 


При $N=2$ и некоторой редукции уравнение (0.1) 


совпадает с известным матричным нелинейным уравнением Шредингера, а 


при $N=3$ --- с матричным модифицированным уравнением 


Кортевега~де~Фриза. Семейство уравнений (0.1) образует так 


называемую иерархию. В статье доказывается, что уравнения (0.1) 


являются дифференциальными, выводятся рекуррентные формулы 


для построения $F_N(x,t,z)$. На этом пути находится бесконечный набор 


полиномиальных законов сохранения для уравнений (0.1). Для 


скалярного случая ($m=1$) близкие результаты известны 


(см. \cite{1}, \cite{2}). Далее используются полученные результаты для 


построения асимптотики матрицы-функции Вейля--Титчмарша $v_N(t,z)$ при 


$z\to\infty$


канонической системы дифференциальных уравнений вида \cite{3}, \cite{4} 


\begin{equation} 


\myfrac{W}{x}=G(x,t,z)W.\tag{0.2} 


\end{equation} 


Этот результат является новым и в скалярном случае. В статье 


введены классы регулярности $\cal P_T$ и выводится теорема существования 


и единственности решения уравнения (0.1) в классе $\cal P_T$ в области 


$0\le x<\infty$, $0\le t\le T$. При этом предполагаются известными 


только начальные данные. 


 


Интересно отметить, что для стационарной линейной задачи 


\begin{equation} 


\frac{d W}{d x}=G(x,0,z)W\tag{0.3} 


\end{equation} 


также находится асимптотика соответствующей функции Вейля--Титчмарша 


$v(z)$. При этом используются шаг за шагом нестационарные  


нелинейные уравнения иерархии. 


 


Таким образом, в статье выявлена связь уравнений иерархии, 


законов сохранения и коэффициентов асимптотики функции 


Вейля--Титчмарша. 


 


\section{Построение иерархии уравнений} 


 


1. Пусть матрицы-функции $G(x,t,z)$ и $F_N(x,t,z)$ имеют вид 


\begin{align} 


G(x,t,z)&=g_1(x,t)+z g_0(x,t),\\ 


F_N(x,t,z)&=f_N(x,t)+z f_{N-1}(x,t)+\dotsb+z^N f_0(x,t). 


\end{align} 


Матрицы-функции $g_k(x,t)$ ($k=0,1$) и $f_k(x,t)$ ($0\le k\le N$) 


порядка $2m\times 2m$ 


предполагаются дифференцируемыми. Кроме того, будем 


далее считать 


\begin{gather} 


g_0=f_0=ij,\quad 


j=\begin{bmatrix} E_m & 0 \\ 0 & -E_m \end{bmatrix}, \\ 


g_1(x,t)=\begin{bmatrix} 0 & R_{12}(x,t) \\ R_{21}(x,t) & 0 


\end{bmatrix}, 


\end{gather} 


где $R_{12}(x,t)$, $R_{21}(x,t)$ --- матрицы порядка $m\times m$. 


 


Рассмотрим матричные уравнения 


\begin{align} 


\myfrac{W(x,t,z)}{x}&=G(x,t,z)W(x,t,z),\\ 


\myfrac{W(x,t,z)}{t}&=F_N(x,t,z)W(x,t,z). 


\end{align} 


Как известно, соотношение (0.1) является условием совместности 


уравнений (1.5) и (1.6). 


 


В данном параграфе выведем рекуррентные формулы для построения 


$f_k(x,t)$ ($k>0$), т.\,е.~дадим метод построения 


иерархии (0.1). Из (0.1) и (1.1)--(1.4) следует 


\begin{align} 


-\myfrac{f_k(x,t)}{x}+[g_1,f_k]+[g_0,f_{k+1}]&=0,\quad 


0\le k<N,\\ 


\myfrac{g_1}{t}-\myfrac{f_N}{x}+[g_1,f_N]&=0. 


\end{align} 


Отметим, что соотношения (1.7) не зависят от $N$. 


 


2. Существенную роль в дальнейшем играет матричное нелинейное  


уравнение 


\begin{equation} 


\myfrac{}{x}w(x,t,z)=2i z w(x,t,z)-w(x,t,z) 


R_{21}(x,t)w(x,t,z)+R_{12}(x,t), 


\end{equation} 


где $w(x,t,z)$ --- матрица-функция порядка $m\times m$. 


 


Из общей теории асимптотических рядов следует утверждение \cite{5}. 


 


\begin{thms} Пусть матрицы $R_{12}(x,t)$, $R_{21}(x,t)$ 


аналитичны по $x$ в области, содержащей полуось $x\ge 0$. Тогда 


существует аналитическое по $x$ $(x\ge 0)$ и $z$ $(|z|>\Delta_0)$ 


решение $w(x,t,z)$, допускающее разложение 


\begin{equation} 


w(x,t,z)=\sum^\infty_{k=1}w_k(x,t,z)/(i z)^k 


\end{equation} 


$(\Delta_0$ зависит от $x$ и $t)$. 


\end{thms} 


 


При этом справедливы равенства 


\begin{equation} 


%%% \left\{ 


\begin{split} 


w_1(x,t)&=-\frac{1}{2}R_{12}(x,t),\\ 


w_{k+1}(x,t)&=\frac{1}{2}\bigg[\myfrac{}{x}w_k(x,t)+ 


\sum_{p+s=k}w_p(x,t)R_{21}(x,t)w_s(x,t)\bigg], 


\end{split} 


\quad p,s\ge 1.


%%% \right. 


\end{equation} 


 


Пусть матрица-функция $\varphi(x,t,z)$ порядка $m\times m$ является 


решением уравнения 


\begin{equation} 


\myfrac{\varphi}{x}=R_{12}w\varphi,\quad 


\varphi(0,t,z)=E_m. 


\end{equation} 


При помощи $w$ и $\varphi$ введем матрицы-функции 


\begin{equation} 


\mu=R_{21}w,\quad \varphi_1=e^{-izx}w\varphi,\quad 


\varphi_2=e^{-izx}\varphi. 


\end{equation} 


Из (1.8), (1.12), (1.13) следует, что матрицы $\varphi_1$ и $\varphi_2$ 


удовлетворяют системе уравнений 


\begin{equation} 


%%% \left\{ 


\begin{split} 


\myfrac{\varphi_1}{x}&=iz\varphi_1+R_{12}\varphi_2,\\ 


\myfrac{\varphi_2}{x}&=-iz\varphi_2+R_{21}\varphi_1. 


\end{split} 


%%% \right. 


\end{equation} 


 


3. Запишем уравнение, аналогичное уравнению (1.9),


\begin{equation} 


\myfrac{}{x}\wt w(x,t,z)=-2iz\wt w(x,t,z)+R_{21}(x,t)- 


\wt w(x,t,z)R_{12}(x,t)\wt w(x,t,z). 


\end{equation} 


Из теоремы 1.1 следует, что решение (1.15) допускает разложение 


\begin{equation} 


\wt w(x,t,z)=\sum^\infty_{k=1}\wt w_k(x,t)/(-iz)^k. 


\end{equation} 


Из (1.16) выводим следующие аналоги соотношений (1.11): 


\begin{align} 


\wt w_1&=-\frac{1}{2}R_{21},\\ 


\wt w_{k+1}&=\frac{1}{2}\bigg(\myfrac{\wt w_k}{x}+ 


\sum_{p+s=k}\wt w_p R_{12}\wt w_s\bigg);\quad p,s,k\ge 1. 


\end{align} 


Матрица $\wt\varphi(x,t,z)$ является решением уравнения 


\begin{equation} 


\myfrac{\wt\varphi}{x}=R_{12}\wt w\wt\varphi,\quad 


\wt\varphi(0,t,z)=E_m. 


\end{equation} 


Введем теперь матрицы-функции 


\begin{equation} 


\wt\mu=R_{12}\wt w,\quad \wt\varphi_1=e^{ixz}\wt\varphi, 


\quad \wt\varphi_2=e^{ixz}\wt w\wt\varphi. 


\end{equation} 


Из (1.15) и (1.19), (1.20) следует, что $\wt\varphi_1(x,t,z)$ и 


$\wt\varphi_2(x,t,z)$ 


также удовлетворяют системе (1.14). Таким образом, доказана 


 


\begin{thms} Пусть выполнены условия теоремы $1.1$. Тогда матрица 


\begin{equation} 


W(x,t,z)= 


\begin{bmatrix} 


\wt\varphi_1(x,t,z)& \varphi_1(x,t,z) \\ 


\wt\varphi_2(x,t,z) & \varphi_2(x,t,z) 


\end{bmatrix} 


\end{equation} 


удовлетворяет системе $(1.5)$. 


\end{thms} 


 


Заметим, что в силу (1.12), (1.19) и разложений (1.9), (1.16) 


матрица $W(x,t,z)$ обратима в окрестности $z=\infty$. 


 


\begin{cors} Матрица-функция 


\begin{equation} 


F(x,t,z)=W(x,t,z)(ij)W^{-1}(x,t,z) 


\end{equation} 


удовлетворяет следующему соотношению: 


\begin{equation} 


\myfrac{F}{x}=[G,F]. 


\end{equation} 


\end{cors} 


 


Из (1.13) и (1.20) вытекают соотношения 


\begin{equation} 


\varphi_1\cdot\varphi^{-1}_2=w,\quad \wt\varphi_2\cdot 


\wt\varphi^{-1}_1=\wt w. 


\end{equation} 


В силу (1.21)--(1.24) равенство (1.22) может быть записано 


в виде 


\begin{equation} 


F(x,t,z)=W_1(x,t,z)(ij)W^{-1}_1(x,t,z), 


\end{equation} 


где 


$$ 


W_1(x,t,z)= 


\begin{bmatrix} 


E_m & w(x,t,z) \\ 


\wt w(x,t,z) & E_m 


\end{bmatrix} 


. 


$$ 


 


\begin{cors} Матрица-функция $F(x,t,z)$ в окрестности 


$z=\infty$ допускает разложение 


\begin{equation} 


F(x,t,z)=ij+\sum^\infty_{k=1}f_k(x,t)/z^k. 


\end{equation} 


\end{cors} 


 


Переписывая (1.25) в виде 


$$ 


F(x,t,z)W_1(x,t,z)=W_1(x,t,z)ij 


$$ 


и пользуясь асимптотическими разложениями (1.9), (1.16), (1.26), 


выводим следующую рекуррентную формулу: 


\begin{multline} 


f_p(x,t)=-\sum^{p-1}_{k=0}f_k(x,t) 


\begin{bmatrix} 


0 & w_{p-k}(x,t)/i^{p-k}\\ \wt w_{p-k}(x,t)/(-i)^{p-k} & 0 


\end{bmatrix} 


- \\ -


\begin{bmatrix} 


0 & w_p(x,t)/i^{p-1} \\ \wt w_p(x,t)/(-i)^{p-1} & 0 


\end{bmatrix} 


,\quad p\ge 1,\ \; f_0=ij. 


\end{multline} 


Из соотношений (1.1), (1.23), (1.26) выводим, что коэффициенты 


$f_k(x,t)$ разложения (1.26) удовлетворяют соотношениям (1.7) 


при всех $k\ge 0$. 


 


Запишем $f_p(x,t)$ в блочном виде 


$$ 


f_p(x,t)=\{f_{kl}(x,t,p)\}^2_{kl=1}, 


$$ 


где $f_{kl}(x,t,p)$ --- матрицы порядка $m\times m$. 


 


\begin{cors} Пусть коэффициенты $f_p(x,t)$ функции 


$F_N(x,t,z)$ (см. (1.2)\,) определены при помощи формул (1.27). 


Тогда уравнение (0.1) эквивалентно системе 


\begin{equation} 


%%% \left\{ 


\begin{split} 


\myfrac{R_{12}}{t}&=2i f_{12}(x,t,N+1),\\ 


\myfrac{R_{21}}{t}&=-2i f_{21}(x,t,N+1). 


\end{split} 


%%% \right.


\end{equation} 


\end{cors} 


 


Действительно, уравнение (0.1) эквивалентно системе (1.8). 


Учитывая формулу (1.4) и равенство (1.7) при $k=N$, приходим 


к системе (1.28). 


 


\begin{notes} Предложенный выбор коэффициентов $f_p(x,t)$ 


функции $F_N(x,t,z)$ не является единственным. Для нас важно, 


что при таком выборе интересные в прикладном и теоретическом плане 


уравнения оказываются членами построенной иерархии. 


\end{notes} 


 


Из формул (1.11), (1.18) следует 


\begin{gather} 


w_2=-\frac{1}{4}\myfrac{R_{12}}{x},\quad 


\wt w_2=-\frac{1}{4}\myfrac{R_{21}}{x},\\ 


\begin{align} 


w_3&=-\frac{1}{8}\myfrac{^2R_{12}}{x^2}+\frac{1}{8} 


R_{12}R_{21}R_{12},\\ 


\wt w_3&=-\frac{1}{8}\myfrac{^2 R_{21}}{x^2}+ 


\frac{1}{8}R_{21}R_{12}R_{21}. 


\end{align} 


\end{gather} 


Из (1.10), (1.17), (1.29)--(1.31) и (1.27) вытекают соотношения 


\begin{gather} 


f_1(x,t)=\begin{bmatrix} 


0 & R_{12}(x,t) \\ R_{21}(x,t) & 0 \end{bmatrix},\qquad 


f_2(x,t)=\frac{1}{2}i\begin{bmatrix} 


R_{12}R_{21} & -\myfrac{R_{12}}{x} \\ 


\myfrac{R_{21}}{x} & -R_{21}R_{12} 


\end{bmatrix} 


, \\ 


\begin{align} 


f_{12}(x,t,3)&=-\frac{1}{4}\bigg(\myfrac{^2R_{12}}{x^2}- 


2R_{12}R_{21}R_{12}\bigg), \\ 


f_{21}(x,t,3)&=-\frac{1}{4}\bigg(\myfrac{^2R_{21}}{x^2}- 


2R_{21}R_{12}R_{21}\bigg). 


\end{align} 


\end{gather} 


В силу (1.28) и (1.32)--(1.34) уравнения иерархии имеют следующий 


вид: 


 


если $N=1$, то 


$$ 


\myfrac{R_{12}}{t}=\myfrac{R_{12}}{x},\quad 


\myfrac{R_{21}}{t}=\myfrac{R_{21}}{x}; 


$$ 





если $N=2$, то 


\begin{equation} 


%%% \left\{ 


\begin{split} 


\myfrac{R_{12}}{t}&=-\frac{1}{2}i\bigg( 


\myfrac{^2R_{12}}{x^2}-2R_{12}R_{21}R_{12}\bigg),\\ 


\myfrac{R_{21}}{t}&=\frac{1}{2}i\bigg(\myfrac{^2R_{21}}{x^2}- 


2R_{21}R_{12}R_{21}\bigg); 


\end{split} 


%%% \right. 


\end{equation} 


 


если выполнено одно из условий редукции 


\begin{equation} 


R_{12}=\pm R^*_{21}, 


\end{equation} 


то система (1.35) эквивалентна матричному нелинейному уравнению 


Шредингера 


$$ 


\myfrac{R_{12}}{t}=-\frac{1}{2}i\bigg( 


\myfrac{^2R_{12}}{x^2}\mp 2R_{12}R^*_{12}R_{12}\bigg); 


$$ 


 


если выполнено условие редукции (1.36) и $N=3$, то система (1.28) 


эквивалентна матричному модифицированному уравнению  


Кортевега~де~Фриза 


$$ 


\myfrac{R_{12}}{t}=-\frac{1}{4}\bigg( 


\myfrac{^3R_{12}}{x^3}\mp\myfrac{R_{12}}{x} 


R^*_{12}R_{12}\mp 3R_{12}R^*_{12} 


\myfrac{R_{12}}{x}\bigg). 


$$ 


 


\section{Законы сохранения} 


 


1. Законы сохранения для уравнения (1.28) имеют вид 


\begin{equation} 


\myfrac{}{t}\Tr[T(x,t)]+\myfrac{}{x} 


\Tr[S(x,t)]=0, 


\end{equation} 


где $\Tr T$ является плотностью, а $\Tr S$ --- потоком. 


 


2. Из (0.1) и (0.2) следует равенство \cite{5} 


\begin{equation} 


\myfrac{}{t}W(x,t,z)-F_N(x,t,z)W(x,t,z)=W(x,t,z)\Phi(t,z),


\end{equation} 


где $\Phi(t,z)$ --- некоторая матрица порядка $2m\times 2m$. Пусть 


$F_{kl}$ 


и $\Phi_{kl}$ --- блоки матриц $F_N$ и $\Phi$ размерности $m\times m$. 


Из соотношений (1.21) и (2.2) следует 


\begin{equation} 


\myfrac{\varphi_2}{t}=F_{21}\varphi_1+F_{22}\varphi_2 


+\wt\varphi_2\Phi_{12}+\varphi_2\Phi_{22}. 


\end{equation} 


В силу (1.13) и (2.3) имеем 


\begin{equation} 


\myfrac{\varphi}{t}\varphi^{-1}=F_{21}w+F_{22} 


+\wt\varphi_2\Phi_{12}\varphi^{-1}_2+\varphi_2\Phi_{22} 


\varphi^{-1}_2. 


\end{equation} 


Заметим, что верно равенство 


\begin{equation} 


\Tr(\wt\varphi_2\Phi_{12}\varphi^{-1}_2+ 


\varphi_2\Phi_{22}\varphi^{-1}_2)= 


\Tr(\varphi^{-1}_2\wt\varphi_2\Phi_{12}+\Phi_{22}). 


\end{equation} 


Согласно (1.12), (1.13) и (1.19), (1.20) справедливо равенство 


$$ 


\varphi^{-1}_2\wt\varphi_2=e^{2izx}\sum^\infty_{k=0}\alpha_k 


(x,t)/z^k. 


$$ 


Теперь, дифференцируя 


обе части (2.4) по $x$ и учитывая (2.5), 


можем записать 


\begin{equation} 


\Tr\bigg(\myfrac{^2\varphi}{x\partial t}\varphi^{-1}- 


\myfrac{\varphi}{t}\varphi^{-1}\myfrac{\varphi}{x}\varphi^{-1}\bigg)= 


\myfrac{}{x}\Tr(F_{21}w+F_{22})+\myfrac{}{x} 


\Tr(\varphi^{-1}_2\wt\varphi_2\Phi_{12}). 


\end{equation} 


Здесь учитывается, что $\Phi_{22}(t,z)$ не зависит от $x$. 


 


Для вывода законов сохранения нам понадобится следующее 


утверждение. 


 


\begin{thms} 


Пусть в дополнение к условиям теоремы $1.1$ выполняется следующее 


условие$:$ $\myfrac{}{t}R_{12}(x,t)$, $\myfrac{}{t}R_{21}(x,t)$ 


аналитичны по $x$ в области, содержащей полуось $x\ge 0$. Тогда верны 


разложения 


\begin{align} 


\myfrac{w(x,t,z)}{t}&=\sum^\infty_{k=1} 


\myfrac{w_k(x,t)}{t}\Big/ (iz)^k,\quad |z|>\Delta_0,\\ 


\myfrac{\wt w(x,t,z)}{t}&=\sum^\infty_{k=1} 


\myfrac{\wt w_k(x,t)}{t}\Big/ (-iz)^k,\quad |z|>\Delta_0. 


\end{align} 


\end{thms} 


 


\begin{pf} Дифференцируя обе части (1.9) по переменной $t$, 


получаем уравнение относительно $\myfrac{w}{t}$. Применяя общую теорему 


(\cite{5}, с.\,255) 


и равенство (1.10), выводим (2.7). Аналогично выводится (2.8). 


\end{pf} 


 


\begin{lems} Существует функция $c(x,t)$ такая, что 


$$ 


\|\varphi^{\pm 1}(x,t,z)\|+\|\wt\varphi(x,t,z)\|\le c(x,t). 


$$ 


\end{lems} 


 


Доказательство следует непосредственно из равенств (1.10), 


(1.12) и (1.16), (1.19). 


 


\begin{lems} Справедлива оценка 


\begin{equation} 


\bigg\|\myfrac{}{x}(\varphi^{-1}_2\wt\varphi_2\Phi_{12}) 


\bigg\|\le[A_1(x,t)|z|^N+A_2(x,t)] 


e^{-2(\im z)x},\quad \im z\ge 0. 


\end{equation} 


\end{lems} 


 


\begin{pf} В силу (1.13), (1.20) верно равенство 


$$ 


\varphi^{-1}_2\wt\varphi_2=e^{2izx}\varphi^{-1}\wt w\wt\varphi. 


$$ 


Из леммы 2.1 следует 


\begin{equation} 


\|\varphi^{-1}_2\wt\varphi_2\|\le c(x,t) 


e^{-2(\im z)x}/|z|,\quad \im z\ge 0. 


\end{equation} 


Учитывая (1.12), (1.19) и (2.10), получаем 


\begin{equation} 


\bigg\|\myfrac{}{x}(\varphi^{-1}_2\wt\varphi_2)\bigg\|\le 


c(x,t)e^{-2(\im z)x},\quad \im z\ge 0. 


\end{equation} 


Рассмотрим теперь $\Phi(t,z)$. Согласно (2.2) имеем 


$$ 


\Phi(t,z)=W(0,t,z)^{-1}\bigg[\myfrac{}{t} 


W(0,t,z)-F_N(0,t,z)W(0,t,z)\bigg], 


$$ 


т.\,е. 


\begin{equation} 


\|\Phi(t,z)\|\le c_1(t)+c_2(t)|z|^N,\quad \im z\ge 0. 


\end{equation} 


Из соотношений (2.9), (2.12) вытекает доказываемая оценка (2.9). 


 


Воспользуемся равенством 


\begin{equation} 


\myfrac{}{t}\bigg(\myfrac{\varphi}{x}\varphi^{-1}\bigg)= 


\myfrac{^2\varphi}{t\partial x}\varphi^{-1}- 


\myfrac{\varphi}{x}\varphi^{-1}\myfrac{\varphi}{t}\varphi^{-1}. 


\end{equation} 


Так как 


$$ 


\Tr\bigg(\myfrac{\varphi}{x}\varphi^{-1} 


\myfrac{\varphi}{t}\varphi^{-1}\bigg)=\Tr\bigg( 


\myfrac{\varphi}{t}\varphi^{-1} 


\myfrac{\varphi}{x}\varphi^{-1}\bigg), 


$$ 


то из (1.12), (2.6), (2.13) и леммы 2.2 получаем равенство 


\begin{equation} 


\myfrac{}{t}\Tr\mu=\myfrac{}{x}\Tr(F_{21}w+F_{22})+ 


O(e^{-2(\im z)x}|z|^N). 


\end{equation} 


 


Перепишем равенство (1.2) в виде 


$$ 


F_N(x,t,z)=[f_0(x,t)+\tfrac{1}{z}f_1(x,t)+\dotsb+ 


\tfrac{1}{z^N}f_N(x,t)]z^N. 


$$ 


Положим 


\begin{equation} 


\mu_n(x,t)=R_{21}(x,t)w_n(x,t). 


\end{equation} 


Приравнивая коэффициенты при одинаковых степенях $z$ в обеих частях 


(2.14), выводим утверждение. 


\end{pf} 


 


\begin{thms} Для уравнений иерархии $(1.28)$ справедливы следующие 


законы сохранения$:$ 


\begin{equation} 


\myfrac{}{t}\Tr\mu_n(x,t)=\myfrac{}{x}i^n\Tr 


\sum_{m+s=n+N}f_{21}(x,t,m)w_s(x,t)/i^s, 


\end{equation} 


где $0\le m\le N$, $n\ge 1$, $s\ge 1$. 


\end{thms} 


 


\begin{notes} Плотности $\Tr\mu_n$ не зависят от $N$, т.\,е.~являются 


общими 


для всей иерархии (1.28). 


\end{notes} 


 


\begin{notes} Пользуясь формулами (1.10), (1.29), (1.30) и 


(2.15), имеем 


$$


\mu_1=-\frac{1}{2}R_{21}R_{12},\quad 


\mu_2=-\frac{1}{4}R_{21}\myfrac{R_{12}}{x},\quad 


\mu_3=-\frac{1}{8}R_{21}\bigg(\myfrac{^2R_{12}}{x^2}- 


R_{12}R_{21}R_{12}\bigg). 


$$


\end{notes} 


 


\begin{notes} Если выполнено условие $R_{12}=\pm R^*_{21}$, то 


$\mu_1=\mp R^*_{12}R_{12}$, т.\,е.~$\mu_1$ является знакоопределенной 


матрицей. 


Значит, соответствующая плотность может трактоваться как энергия. 


\end{notes} 


 


Отметим еще, что в этом важном частном случае ($n=1$) формула (2.16) 


упрощается и принимает вид 


\begin{equation} 


\myfrac{}{t}\Tr\mu_1(x,t)=\myfrac{}{x}(-i) 


\Tr f_{22}(x,t,N+1). 


\end{equation} 


Формула (2.17) вытекает непосредственно из (2.16), если 


учесть, что в силу (2.14) верно равенство 


$$ 


f_{22}(x,t,N+1)=-\sum^N_{k=0}f_{21}(x,t,k) 


w_{N+1-k}(x,t)/i^{N+1-k}. 


$$ 


 


\section{Асимптотика функции Вейля--Титчмарша} 


 


1. В этом параграфе отдельно рассмотрим случай, когда  


выполнено условие редукции 


\begin{equation} 


R_{12}(x,t)=R^*_{21}(x,t). 


\end{equation} 


В этом случае линейная задача (0.2) является самосопряженной, 


т.\,е.~могут быть использованы результаты спектральной теории 


\cite{1}. В 


частности, системе (0.2) соответствует матрица-функция Вейля--Титчмарша 


$v(t,z)$ порядка $m\times m$ которая определяется при 


помощи следующего неравенства 


$$ 


\int^\infty_0[E_m,i v^*(t,z)]U w^*(x,t,z)w(x,t,z) 


U^*\fracwithdelims[][0pt]{E_m}{-i v(t,z)}dx<\infty, 


$$ 


где $w(x,t,z)$ --- решение системы (0.2), удовлетворяющее условию 


$$ 


w(0,t,z)=E_{2m},


$$ 


а 


\begin{equation} 


U=\frac{1}{\sqrt{2}} 


\begin{bmatrix} 


E_m & -E_m \\ E_m & \phantom{-}E_m 


\end{bmatrix} 


. 


\end{equation} 


Нам понадобится следующая 


 


\begin{thms}[{\series{m}\shape{n}\selectfont\cite{2}}] Пусть решение 


системы $(1.28)$ удовлетворяет условию редукции $(3.1)$ и неравенству 


$$ 


\bigg\|\myfrac{^k R_{12}(x,t)}{x^k}\bigg\|\le M,\quad 


0\le k\le N,\quad 0\le x<\infty,\quad 0\le t\le T. 


$$ 


Тогда соответствующая система $(0.2)$ имеет единственную функцию 


Вейля--Титчмарша $v_N(t,z)$, которая определяется при помощи 


равенства 


\begin{equation} 


v_N(t,z)=i\{r_{11}(t,z,N)[-iv(z)]+r_{12}(t,z,N)\} 


\{r_{21}(t,z,N)[-iv(z)]+r_{22}(t,z,N)\}^{-1}, 


\end{equation} 


причем $v(z)=v(0,z)$, а матрица 


$r_N(t,z)=\{r_{kl}(t,z,N)\}^2_{kl=1}$ 


определяется соотношениями 


\begin{equation} 


\frac{d r_N(t,z)}{dt}=-\wt F^*_N(0,t,\ov z) 


r_N(t,z),\quad r_N(0,z)=E_{2m}, 


\end{equation} 


где 


\begin{equation} 


\wt F_N(0,t,z)=U F_N(0,t,z)U^*. 


\end{equation} 


\end{thms} 


 


Запишем матрицу $\wt F_N(0,t,z)$ в блочном виде 


$$ 


\wt F_N(0,t,z)=\{\wt F_{kl}(t,z,N)\}^2_{kl=1}, 


$$ 


где все блоки имеют размерности $m\times m$. Дифференцируя по $t$ обе 


части (3.3) и учитывая (3.4), получаем 


\begin{equation} 


\frac{d v_N}{d t}=-\wt F^*_{11}(t,\ov z,N)v_N-i 


\wt F^*_{21}(t,\ov z,N)+ v_N\wt F^*_{22}(t,\ov z,N)-i 


v_N\wt F^*_{12}(t,\ov z,N)v_N. 


\end{equation} 


 


2. Далее будем предполагать, что $v_N(t,z)$ в окрестности 


$z=\infty$ допускает разложение 


\begin{equation} 


v_N(t,z)=\sum^\infty_{k=0}\alpha_k(t,N)/z^k,\quad 


\alpha_0(t,N)=i E_m. 


\end{equation} 


 


\begin{defns} Будем говорить, что решение $R_{12}(x,t)$  


системы (1.28) принадлежит классу регулярности $\cal P_T$, если 


функция Вейля--Титчмарша $v_N(t,z)$ соответствующей системы (0.2) 


допускает в 


окрестности $z=\infty$ при $0\le t\le T$ представление (3.7). 


\end{defns} 


 


В силу (1.2) и (3.4) матрица-функция $\wt F^*_{kl}(t,\ov z,N)$ 


допускает представление 


$$ 


\wt F^*_{kl}(t,\ov z,N)=\sum^N_{s=0}\psi_{kl}(t,s)z^{N-s}.


$$ 


Подставляя формулу (3.7) в уравнение (3.6) и приравнивая 


коэффициенты при одинаковых степенях $z$, получаем 


$$ 


\frac{d\alpha_k(t,N)}{dt}=\sum_{p+s-N=k} 


\psi_{11}(t,s)\alpha_p(t,N)+\sum_{p+s-N=k} 


\alpha_p(t,N)\psi_{22}(t,s)-i 


\sum_{p+l+s-N=k}\alpha_p(t,N)\psi_{12}(t,s) 


\alpha_l(t,N), 


$$ 


где $p\ge 0$, $l\ge 0$, $0\le s\le N$, $k\ge 1$. 


 


Кроме того, из соотношений (3.6), (3.7) следует 


\begin{multline} 


-\sum_{p+s=m}\psi_{11}(t,s)\alpha_p(t,N)+ 


\sum_{p+s=m}\alpha_p(t,N)\psi_{22}(t,s)- \\


-\sum_{p+l+s=m}\alpha_p(t,N)\psi_{12}(t,s)\alpha_l(t,N)- 


i\psi_{21}(t,m)=0, 


\end{multline} 


где $p\ge 0$, $l\ge 0$, $0\le s\le m\le N$. 


 


Согласно (1.3), (3.2) и (3.5) верно равенство 


\begin{equation} 


\begin{bmatrix} 


\psi_{11}(t,0) & \psi_{12}(t,0) \\ 


\psi_{21}(t,0) & \psi_{22}(t,0) 


\end{bmatrix} 


=i 


\begin{bmatrix} 


0 & E_m \\ E_m & 0 


\end{bmatrix} 


. 


\end{equation} 


Учитывая (3.9), перепишем (3.8) в виде 


\begin{multline} 


2i\alpha_m(t,N)=\sum_{p+s=m}\psi_{11}(t,s)\alpha_p(t,N)- 


\sum_{p+s=m}\alpha_p(t,N)\psi_{22}(t,s)+i\psi_{21}(t,m)+ \\


+i\sum_{p+l+s=m}\alpha_p(t,N)\psi_{12}(t,s)\alpha_l(t,N), 


\end{multline} 


где $0\le s\le m\le N$, $0\le p<m$, $0\le l<m$. 


Так как коэффициенты $\psi_{ij}(t,s)$ рекуррентного соотношения 


(3.10) не зависят от $N$, то верно 


 


\begin{stats} 


Пусть $N_1\le N_2$, тогда справедливы равенства 


$$ 


\alpha_m(t,N_1)=\alpha_m(t,N_2),\quad 0\le m\le N_1. 


$$ 


\end{stats} 


 


\begin{notes} Из формул (1.2), (1.27) и (3.5), (3.10) 


вытекают, в частности, равенства 


\begin{align*} 


\alpha_1(t,N)&=-R^*_{12}(0,t),\quad N\ge 1,\\ 


\alpha_2(t,N)&=-\frac{1}{2}i\bigg[\myfrac{}{x} 


R^*_{12}(x,t)\mid_{x=0}+2R^*_{12}(0,t)^2\bigg],\quad N\ge 2. 


\end{align*} 


\end{notes} 


 


3. Как и в скалярном случае \cite{6}, функции $\wt F^*_{kl}(x,t,N)$ 


могут быть выражены через $\alpha_m(t,N)$, при этом  


следует пользоваться формулами (1.2), (1.27) и (3.5), (3.10).


Опустим соответствующие выражения из-за их громоздкости. После этого 


уравнение (3.6) связывает только коэффициенты $\alpha_m(t,N)$ и не  


содержит $R_{12}(0,t)$, 


$\myfrac{^k}{x^k}R_{12}(x,t)\mid_{x=0}$. Повторяя теперь рассуждения 


статьи \cite{6}, выводим 


 


\begin{stats} Пусть $\im v_0(z)>0$ при $\im z\ge 0$ и 


$$ 


v_0(z)=i E_m+\sum^\infty_{k=1}\alpha_k(0)/z^k,\quad 


|z|\ge |z_0|>0. 


$$ 


Тогда уравнение (3.6) имеет одно и только одно решение $v_N(t,z)$ 


такое, что при $0\le t\le T_N$ и $|z|\ge|z_N|$ верно представление 


(3.7), причем 


$$ 


v_N(0,z)=v_0(z),\quad \im v_N(t,z)\ge 0,\quad \im z\ge 0. 


$$ 


\end{stats} 


 


Из утверждения 3.2 и результатов статьи \cite{1} непосредственно 


следует 


 


\begin{thms} Пусть матрица-функция 


\begin{equation} 


R_{12}(x,0)=\Phi(z) 


\end{equation} 


такова, что матрица-функция Вейля--Титчмарша $v_0(z)$ соответствующей 


системы $(0.3)$ удовлетворяет условиям утверждения $3.2$. 


Тогда при некотором $T_N>0$ система $(1.28)$, $(3.1)$ имеет одно и 


только одно решение $R_{12}(x,t,N)$ $(0\le x<\infty$, 


$0\le t\le T_N)$, принадлежащее 


классу $\cal P_{T_N}$ и удовлетворяющее условию $(3.11)$.


\end{thms} 


 


В частности, теорема 3.2 применима к нелинейному уравнению  


Шредингера ($N=2$) и модифицированному уравнению 


Кортевега~де~Фриза ($N=3$). 


 


\begin{notes} Из теоремы 3.2 следует, что для однозначной 


разрешимости уравнений иерархии (1.28), (3.1) при $0\le x<\infty$, 


$0\le t\le T$ в классе $\cal P_T$ достаточно начальных данных и не 


требуется присоединение граничных данных. 


\end{notes} 


 


4. Формулы (3.10) верны и при $t=0$. В этом случае получаем 


равенство 


\begin{multline} 


2i\alpha_m(0)=\sum_{p+s=m}\psi_{11}(0,s)\alpha_p(0)- 


\sum_{p+s=m}\alpha_p(0)\psi_{22}(0,s)+\\ 


+i\sum_{p+l+s=m}\alpha_p(0)\psi_{12}(0,s)\alpha_l(0)+ 


i\psi_{21}(0,m), 


\end{multline} 


где $0\le p<m$, $0\le l<m$, $s\ge 0$, $m>0$. 


Формула (3.12) связывает коэффициенты разложения $\alpha_m(0)$  


функции $v_0(z)$ со значениями $\Phi^{(k)}(0)$. 


 


Интересно отметить, что формула (3.12), относящаяся к линейной 


стационарной системе (0.3), выведена в результате исследования иерархии 


нелинейных нестационарных систем. 


 


 


\begin{thebibliography}{9} 


 


\bibitem{1} 


Абловиц М., Сигур Х. {\em Солитоны и метод обратной задачи}. 


-- М.: Мир, 1987. -- 480 с. 


\bibitem{2} 


Захаров В.Е., Шабат А.Б. {\em Точная теория двумерной самофокусировки и 


одномерной автомодуляции волн в нелинейной среде } // 


ЖЭТФ. -- 1971. -- Т.\,61. -- С.\,118--134. 


\bibitem{3} 


Сахнович Л.А. {\em Задачи факторизации и операторные тождества} 


// УМН. -- 1986. -- Т.\,41. -- \No\,1. -- С.\,3--55. 


\bibitem{4} 


Сахнович Л.А. {\em Метод операторных тождеств и задачи анализа} // 


Алгебра и анализ. -- 1993. -- Т.\,5. -- \No\,1. -- C.\,3--80.


\bibitem{5} 


Вазов В. {\em Асимптотические разложения решений обыкновенных 


дифференциальных уравнений}. -- М.: Мир, 1968. -- 464 с. 


\bibitem{6} 


Сахнович Л.А. {\em Интегрируемые нелинейные уравнения на полуоси} // 


Укр. матем. журн. -- 1991. -- T.\,43. -- \No\,11. -- С.\,1578--1584. 


 


\end{thebibliography} 


 


\vskip 2cm 


 


\post{\it Украинская государственная}{\it Поступила} 


\post{\it Академия связи}{19.03.1997} 


 


\label{end} 


\end{document} 


